%% kapa: exam=1 round=37 year=2026 period=2 form=A subject=appraisal_law subject_ko=감정평가관계법규 seq=15 number=15
%% source: Q-Net 감정평가사 기출문제 (cst00309) artlSeq=5258924 file=raw/1008/5258924/2026년 제37회 감정평가사 1차 2교시.pdf
%% conversion: kordoc 4.24.51 -> tools/gen1.py (auto); review=proofread-sonnet; flags=none
\begin{question}{15}
부동산 가격공시에 관한 법령상 표준지공시지가에 관한 설명으로 옳은 것은?
\begin{choices}
  \item 표준지의 도로상황은 표준지공시지가의 공시에 포함되어야 하는 것은 아니다.
  \item 표준지공시지가에 이의가 있는 자는 그 공시일부터 30일 이내에 시ㆍ도지사에게 이의를 신청할 수 있다.
  \item 표준지공시지가는 표준지공시지가 조사ㆍ평가보고서에 따른 조사ㆍ평가액의 기하평균치를 기준으로 한다.
  \item 표준지에 지상권이 설정되어 있을 때에는 그 권리가 존재하는 것으로 보고 표준지공시지가를 평가하여야 한다.
  \item 표준지공시지가 조사ㆍ평가보고서에는 토지 형상 및 지세(地勢)가 포함되어야 한다.
\end{choices}
\end{question}
\answer{5}
